% Template for the Journal of Computational Science and Applications (JCSA)
% Published by the University of Lahore -- ISSN 2710-3102
% Official author guidelines: https://journals.uol.edu.pk/JCSA/Author-Guideline
\documentclass[twoside,twocolumn,10pt]{article}
\usepackage{lipsum}  % For sample/placeholder text
\usepackage{fancyhdr}
\usepackage{geometry}
\usepackage{multicol}
\usepackage{abstract}
\usepackage{graphicx}
\usepackage{caption}
\usepackage{titlesec}
\usepackage[utf8]{inputenc} % For UTF-8 input encoding
\usepackage[T1]{fontenc}    % Better font encoding for PDF output
\usepackage{ragged2e}  % Provides the justifying command
\usepackage{amssymb}  % For additional symbols
\usepackage{parskip}  % Automatically handles paragraph spacing
\usepackage{amsmath}  % For advanced math environments
\usepackage{newtxtext}  % Times New Roman font
\usepackage{booktabs}  % For improved table appearance
\usepackage{array}     % For customizing table columns
\usepackage{hyperref} % For hyperlinks in the document
\usepackage[backend=biber,style=numeric,sorting=none]{biblatex} % For bibliography management


% The starting page number is normally assigned by the editorial office
% once the article is scheduled into a specific issue. Authors can leave
% this at 1; production will update it if required.
\setcounter{page}{1}

% Define custom column types for better alignment and spacing
\newcolumntype{L}[1]{>{\raggedright\arraybackslash}p{#1}} % Left aligned column
\newcolumntype{C}[1]{>{\centering\arraybackslash}p{#1}}   % Center aligned column
\newcolumntype{R}[1]{>{\raggedleft\arraybackslash}p{#1}}  % Right aligned column

\usepackage[font=footnotesize, labelfont=bf, justification=raggedright, format=plain]{caption}

\captionsetup[table]{
    labelsep=period,               % Period after label for tables
    justification=raggedright      % Ensure left alignment of captions
}

\captionsetup[figure]{
    labelsep=period,               % Period after label for figures
    justification=raggedright      % Ensure left alignment of captions
}

\addbibresource{sources.bib} %Import the bibliography file
% Customizing bibliography heading
\defbibheading{bibliography}[\refname]{%
  \section{\MakeUppercase{REFERENCES}}}

% Set page margins
\geometry{
    a4paper,
    left=0.8in,
    right=0.8in,
    top=1in,
    bottom=1in
}

% Increase space between columns
\setlength{\columnsep}{15pt}

% Define the centred abstract heading
\renewcommand{\abstractname}{\hspace{1.3em}\textbf{Abstract}}

\titleformat{\section}{\large\bfseries\uppercase}{\thesection.}{1em}{}

\titleformat{\subsection}{\normalfont\bfseries\flushleft}{\thesubsection.}{1em}{}

\titleformat{\subsubsection}{\normalfont\bfseries\itshape\raggedright}{\thesubsubsection.}{1em}{}

% Set font sizes for abstract
\renewcommand{\abstractnamefont}{\normalfont\fontsize{14}{20}\bfseries}
\renewcommand{\absnamepos}{flushleft}
\renewcommand{\abstracttextfont}{\normalfont\fontsize{11}{13}\selectfont}

% Keywords environment
\newcommand{\keywordsname}{Keywords}
\newenvironment{keywords}
  {\par\noindent\hspace{2em}\textbf{\keywordsname:}}
  {\par}

\hyphenpenalty=10000
\exhyphenpenalty=10000
\sloppy

% Set up header and footer styles
\pagestyle{fancy}
\fancyhf{}  % Clear default header and footer

% Odd page headers -- replace with a short running form of your title
\fancyhead[C]{\small Running Title of the Paper}

% Even page headers -- replace with the first author's surname and year
\fancyhead[CE]{\small Author Surname et al. (Year)}

% Footer: Page numbers only (on all pages except the first one)
\fancyfoot[C]{\thepage}

% First page custom header/footer
\fancypagestyle{firstpage}{
  \fancyhf{}  % Clear header and footer
  \fancyhead[LO]{\small Volume X, Issue Y, Year}  % Left header for first page
  \fancyhead[RO]{\small J. Comput. Sci. Appl.}  % Right header for first page
  \fancyfoot[L]{\footnotesize Received: Month Year; Revised: Month Year; Accepted: Month Year\\\textsuperscript{*}Corresponding Author: Full Name \textless{}email@example.com\textgreater{}}
   \fancyfoot[R]{\footnotesize \copyright\ Year by the Author (s). Published by JCSA,\\\textsuperscript{*} under the CC-BY license }
  \renewcommand{\headrulewidth}{0.4pt}  % Add header line
  \renewcommand{\footrulewidth}{0.4pt}  % Add footer line
}

\title{\fontsize{16}{20}\selectfont\textbf{Stem Cell Growth using Mycelium Framework}}
\author{
    \fontsize{12}{14}\selectfont
    Oliver Linger\textsuperscript{1*} \\
    \fontsize{12}{14}\selectfont
    \textsuperscript{1}Department Name, University Name, City, Country 
}
\date{}

\begin{document}

\twocolumn[
\begin{@twocolumnfalse}
% Change "Research Article" to the appropriate article type if needed
% (e.g. Review Article, Short Communication).
\vspace{-0.5cm}  % Adjust spacing as needed between header and "Research Article"
\raggedright%

{\fontsize{11}{13}\selectfont{Research Article}}
\vspace{-0.8cm}  % Adjust spacing between "Research Article" and title

% Now, insert the title and the author information
\maketitle

\thispagestyle{firstpage} % Apply custom header/footer for first page
\begin{abstract}
\fontsize{11}{13}\selectfont
\vspace{1em}
\justifying%
Mycelium

\vspace{-0.1cm}
\begin{keywords}
Place text here
\end{keywords}

\end{abstract}

\vspace{1cm}
\end{@twocolumnfalse}
]
\section{Introduction}
\fontsize{12}{14}\selectfont
Using myclium scaffold to influence stem cell adaptation. 
Creating the stem cell niche by growing the stem cell in a natural biological scaffold.

\section{Literature Review}
\fontsize{12}{14}\selectfont
In~\cite{rufferi_mycleium_based_wound_healing_2023}, the affects of a biological gause on wound healing is observed.
It concludes that the mycelium biomaterials are a promising approach to alternative accelrative wound healing.
Both tested materials, G. lucidum and P. ostreatus enhance collagen I gene expression.
G. lucidum performs marginally better. The $\beta$-glucans in the mycelia are crucial in prompting the body to synthesize types I and III collagen.

In~\cite{Antinori2021Mycelium} it is demonstrated that the mycelia of Pleurotus ostreatus can be directly used as the scaffold for cell growth.
It had good viability, and morphology compared to the control samples used in~\cite{Antinori2021Mycelium}. Human cells were directly plated onto the un-moddified and non-functionalised mycelial scaffold. 
Worth noting; even after Hyphal cell deactivation, Ganoderma lucidum was not fully compatible for direct plating. 
Its water extract contained concentrations of organic acids (Ganoderic acid V and Oleandrin). 
This was detrimental for HDFa cells survivial

In~\cite{Voog2010NicheStemCells}, the stem cell niche is discussed. It details the effect of the microenvironment on stem cell differenciation.
It provides a methodology for defining the cell types that stem cells can grow into, and the condiitons that are needed for this to occur.
Under homeostatic conditions, there is a symbiotic relationship between the stem cells and their niche. 
What must be further investigated is the response of the stem cells and their niche to tissue injury during the disease progression and aging. 
Further study must be done on genetic factors that influence the stem cell niche. 

\section{Methodology and Data}
\fontsize{12}{14}\selectfont
Place text here

\subsection{Example Subsection}
\fontsize{12}{14}\selectfont
Place text here

\subsection{Example Second Subsection}
\fontsize{12}{14}\selectfont
Place text here

\subsubsection{Example Subsubsection}

Place text here

\section{Results and Discussion}

Place text here

\section{Conclusion and Future Work}

Place text here

\section{Declaration}
\fontsize{11}{13}\selectfont

\textbf{Conflict of Interest: } Place text here

\textbf{Author Contribution:} Place text here

% Example for references section

\printbibliography% Use BibLaTeX for printing bibliography
\end{document}